\documentclass[11pt,a4paper]{article}
\usepackage[portuguese]{babel}
\usepackage[utf8]{inputenc}
\usepackage[T1]{fontenc}
\usepackage{geometry}
\geometry{margin=2.5cm}

\title{Recolha de Informação e Infraestrutura Técnica para Monitorização de Doentes Cardíacos}
\author{}
\date{}

\begin{document}
\maketitle

\section*{Introdução}

O presente trabalho insere-se no desenvolvimento de um sistema de
telemonitorização de saúde (\textit{pHealth}) destinado a doentes cardíacos com
miocardiopatia pós-inflamatória.

Com o objetivo de fundamentar a arquitetura técnica e funcional do sistema,
realizou-se uma consulta especializada junto de profissionais de saúde com
profundo conhecimento clínico e experiência prática no acompanhamento desta
patologia. Esta abordagem permitiu identificar os requisitos clínicos críticos
e os parâmetros fisiológicos mais relevantes para uma monitorização eficaz bem
como entender as necessidades e principais desafios para os utentes.

Especificamente, buscou-se compreender:
\begin{itemize}
    \item O perfil clínico e funcional típico dos doentes com esta patologia
    \item Quais os parâmetros fisiológicos que devem ser monitorizados continuamente ou
          periodicamente;
    \item O estado da arte em termos de aparelhos e tecnologias de monitorização
          disponíveis atualmente no mercado;
    \item Qual a infraestrutura técnica necessária para a recolha, processamento e
          armazenamento desses dados de forma segura e eficiente.
\end{itemize}

\section*{Questões Clínicas e Funcionais Avançadas}

\subsection*{Caracterização Clínica}

A miocardiopatia pós-inflamatória é uma condição complexa causada pela
inflamação do músculo cardíaco (miocárdio), frequentemente com componente
genética associada. Esta patologia resulta em disfunção significativa do
sistema de condução elétrica do coração, levando tipicamente a manifestações
clínicas que requerem monitorização rigorosa.

\subsection*{Importância da Monitorização de Frequência e Ritmo Cardíaco}

Um dos aspetos críticos nestes doentes é a alteração da frequência cardíaca.
Enquanto alguns apresentam taquicardia (frequência elevada) como resposta à
disfunção miocárdica, outros podem desenvolver bradicardia excessiva induzida
pela medicação. Por este motivo, é fundamental monitorizar continuamente tanto
a frequência cardíaca como o ritmo cardíaco, de modo a:
\begin{itemize}
    \item Detetar e controlar episódios de taquicardia, frequentemente indicadores de
          arritmias subjacentes;
    \item Identificar bradicardia excessiva causada pelos fármacos utilizados no
          tratamento;
    \item Reconhecer precocemente padrões anormais que possam comprometer a perfusão
          tecidular.
\end{itemize}

\subsection*{Arritmias Potencialmente Fatais}

A principal preocupação clínica, tanto para doentes como para profissionais de
saúde, é a ocorrência de arritmias potencialmente fatais. As mais relevantes
incluem:

\begin{itemize}
    \item \textbf{Fibrilhação auricular:} arritmia supraventricular que aumenta significativamente o risco de tromboembolismo;

    \item \textbf{Extrasístoles ventriculares:} batimentos cardíacos prematuros originados nos ventrículos. Quando frequentes ou agrupadas, podem evoluir para arritmias mais graves;

    \item \textbf{Fibrilhação ventricular:} uma das situações mais críticas, caracterizada por atividade elétrica caótica e desorganizada nos ventrículos. Nesta condição, 
    a contração ventricular torna-se ineficaz, impedindo a ejeção de sangue adequada. A diástole extremamente curta não permite o enchimento ventricular, resultando em perda 
    imediata da perfusão sistémica. Esta condição manifesta-se frequentemente com palpitações vigorosas e representa uma emergência médica absoluta, requerendo intervenção imediata 
    (ressuscitação cardiopulmonar e desfibrilhação) para restaurar um ritmo cardíaco viável e prevenir morte súbita.
\end{itemize}

Consequentemente, a detecção precoce e o reconhecimento rápido destas arritmias
constituem objetivos prioritários da monitorização contínua.

\section*{Infraestrutura Tecnológica}
TODOf
\end{document}